% Build in this directory: pdflatex -interaction=nonstopmode -halt-on-error class12_handout.tex
\documentclass[a4paper,10pt,pdflatex,ja=standard,jafont=ipaex-type1,
  magstyle=real,scale=1]{bxjsarticle}
\usepackage{lmodern}
\usepackage{amsmath,amssymb,mathtools}
\usepackage{geometry}
\geometry{reset,left=22mm,right=22mm,top=20mm,bottom=17mm,
  headheight=13pt,headsep=6mm,footskip=10mm}
\usepackage{xcolor}
\usepackage{enumitem}
\usepackage{array,booktabs,tabularx}
\usepackage{fancyhdr}

\definecolor{ink}{HTML}{193A5C}
\definecolor{accent}{HTML}{237D82}
\definecolor{soft}{HTML}{EAF3F3}
\definecolor{muted}{HTML}{52616B}
\newcommand{\coursename}{解析学 II}
\newcommand{\handoutnumber}{12}
\pagestyle{fancy}
\fancyhf{}
\fancyhead[L]{\small\color{ink}\coursename{}・授業要約と演習}
\fancyhead[R]{\small\color{ink}第\handoutnumber 回}
\fancyfoot[C]{\small\color{ink}\thepage/2}
\renewcommand{\headrulewidth}{0.3pt}
\setlength{\parindent}{0pt}
\setlength{\parskip}{2pt}
\setlength{\abovedisplayskip}{5pt}
\setlength{\belowdisplayskip}{5pt}
\setlength{\abovedisplayshortskip}{3pt}
\setlength{\belowdisplayshortskip}{4pt}
\setlist[itemize]{leftmargin=1.8em,itemsep=1pt,topsep=2pt,parsep=0pt}

\newcommand{\handouttitle}[1]{%
  \begin{center}
    {\Large\color{ink}#1}\par\vspace{3mm}
    {\color{accent}\rule{0.88\linewidth}{0.7pt}}
  \end{center}\vspace{-2mm}}
\newcommand{\handout}[2]{%
  \renewcommand{\handoutnumber}{#1}%
  \handouttitle{第#1回\quad #2}}
\newcommand{\topic}[1]{\par\vspace{5pt}%
  {\large\sffamily\mdseries\color{accent}#1}\par\vspace{2pt}}
\newcommand{\keybox}[1]{\par\vspace{3pt}\begingroup
  \setlength{\fboxsep}{4pt}%
  \colorbox{soft}{\parbox{\dimexpr\linewidth-2\fboxsep\relax}{\small #1}}%
  \endgroup\par\vspace{2pt}}
\newenvironment{problems}
  {\begin{enumerate}[leftmargin=2em,itemsep=3pt,topsep=3pt,parsep=0pt,
    font=\color{accent}]}
  {\end{enumerate}}


\newcommand{\examplelabel}[1]{{\sffamily\mdseries\color{ink}#1}}
\newcommand{\answerroom}{\par\vspace{16mm}}

\begin{document}
\fontsize{10pt}{13.5pt}\selectfont
\setlength{\abovedisplayskip}{4pt}
\setlength{\belowdisplayskip}{4pt}
\handout{12}{曲線・曲面の向き}
\topic{曲線の向きと線積分の符号}
曲線の表示 $\vec r(t)$（$a\le t\le b$）は、$t$ の増加方向に向きを与える。
逆向きの表示は $\widetilde{\vec r}(t)=\vec r(a+b-t)$ である。
$$
 \int_{-C}\vec F\cdot d\vec r=-\int_C\vec F\cdot d\vec r,
 \qquad \int_{-C}f\,ds=\int_Cf\,ds.
$$
平面の単純閉曲線は、囲む領域を左に見る向きが正である。
開いた線分には自然な「正」の向きはなく、始点・終点で指定する。

\topic{曲面の媒介変数表示と正則性}
$D\subset\mathbb R^2$ 上の $C^1$ 級写像
$\vec r(u,v)=(x(u,v),y(u,v),z(u,v))$ で曲面 $S$ を表す。
内部で $\vec r_u\times\vec r_v\ne\vec0$ なら、2本の接ベクトルは独立である。
$$
 \vec N=\vec r_u\times\vec r_v,\qquad
 \vec n=\frac{\vec N}{|\vec N|},\qquad dS=|\vec N|\,du\,dv.
$$
表示は内部で1対1とする。球面などは必要に応じて複数の座標で覆う。

\topic{曲面の向きと法線}
連続な単位法線 $\vec n$ を選べる曲面を向き付け可能という。
連結な曲面の向きは $\vec n$ と $-\vec n$ の2通り。
パラメータの順序を交換すると、外積の向きが逆になる。
$$
 \vec r_v\times\vec r_u=-\vec r_u\times\vec r_v.
$$
グラフ $z=f(x,y)$ を $(x,y)$ の順で表すと $\vec N=(-f_x,-f_y,1)$ は上向き。
閉曲面が立体を囲む場合は、外向きか内向きかを指定する。

\topic{例：円柱と球面を外向きに表す}
半径1の円柱側面では $\vec r(\theta,z)=(\cos\theta,\sin\theta,z)$ とすれば
$$
 \vec r_\theta\times\vec r_z=(\cos\theta,\sin\theta,0)
$$
は外向きである。
球面では $\vec r(\alpha,\theta)=(\sin\alpha\cos\theta,\sin\alpha\sin\theta,\cos\alpha)$ として
$$
 0<\alpha<\pi,\quad0\le\theta<2\pi,\qquad
 \vec r_\alpha\times\vec r_\theta=\sin\alpha\,\vec r.
$$
これも外向きである。極で外積が0になるのは、この座標表示の退化であり、球面自体の特異点ではない。

\topic{曲面の向きから境界の向きを決める}
右手の親指を法線方向に向けたとき、他の指が曲がる方向を境界の正の向きとする。
曲面上を進み、選んだ法線側から見て内部が左になる向きに対応する。
$$
 \text{円板の法線が }+z\text{ 方向}\ \Longrightarrow\
 \partial S\text{ は }+z\text{ 側から見て反時計回り}.
$$
法線を反転すると、誘導される境界の向きも反転する。
\keybox{媒介変数の順序・外積の向き・指定された法線・境界の進む方向を、一組として確認する。}

\newpage
\handouttitle{第12回\quad 演習問題}
計算過程を記し、必要な条件・積分領域・向きを明示すること。
\begin{problems}
\item $(-1,1)$ から $(0,1)$ への線分と、その逆向きの線分をそれぞれ媒介変数で表せ。$0\le t\le1$ を用いること。\answerroom
\item $a,b>0$ の楕円 $x^2/a^2+y^2/b^2=1$ を、反時計回り・時計回りの両方で表せ。$t=0$ での位置と速度で向きを確認せよ。\answerroom
\item $z=x^2+y^2$、$x^2+y^2\le1$ の曲面を上向きに表し、接ベクトル・単位法線・面積要素を求めよ。\answerroom
\item 半径 $R>0$、高さ $H>0$ の円柱側面を外向きに表し、外積を計算せよ。パラメータの順序を交換すると何が変わるか。\answerroom
\item 単位球面の表示 $\vec r(\alpha,\theta)=(\sin\alpha\cos\theta,\sin\alpha\sin\theta,\cos\alpha)$ の外積を計算し、外向きであることを確かめよ。極をどう扱うかも説明せよ。\answerroom
\item 単位立方体の上面 $z=1$ と下面 $z=0$ を外向きに向き付ける。それぞれの境界を $+z$ 側から見ると、正の向きは時計回りか反時計回りか。理由も述べよ。\answerroom
\end{problems}
\end{document}
